\documentclass[11pt]{beamer}
\usetheme{Madrid}
\usecolortheme{default}

\setbeamertemplate{navigation symbols}{}
\setbeamertemplate{footline}[frame number]
\setbeamercolor{structure}{fg=blue!60!black}
\setbeamercolor{title}{fg=white}
\setbeamercolor{subtitle}{fg=white}
\setbeamercolor{author}{fg=black}
\setbeamercolor{institute}{fg=black}
\setbeamercolor{date}{fg=black}
\setbeamercolor{frametitle}{bg=white,fg=black}

\usepackage{amsmath,amssymb}
\usepackage{booktabs}
\usepackage[round,authoryear]{natbib}
\usepackage{graphicx}
\usepackage{tikz}
\usetikzlibrary{arrows.meta,positioning}

\title{Outcome Weights, Kappa Estimators,\\ and Covariate Balance}
\subtitle{Auditing IV and Causal ML Estimartors through Outcome Weights}
\author{Alexander Unger}
\institute{University of T\"ubingen \\ Supervisor: JProf.\ Michael Knaus}


\begin{document}

% ---------- TITLE ----------
\begin{frame}
\titlepage
\end{frame}

% ---------- OUTLINE ----------
\begin{frame}{Roadmap}
\tableofcontents
\end{frame}

% ============================================================
\section{Problem statement}
% ============================================================




\begin{frame}{The effect on millitary enlisting on earnings}
In causal inference, we often want to measure the impact of a policy or treatment on a later outcome for example, the effect of having served in the military on lifetime wages.

\vspace{0.2cm}
One might expect ambiguous effects \textit{a priori}: military service takes years away from civilian careers, but could also provide skills, discipline, or benefits (e.g., the GI Bill) that raise later earnings.

\vspace{0.3cm}
\citet{sloczynski2025abadie} study exactly this question, using the Vietnam-era draft lottery \citet{angrist1990lifetime} as an instrument for veteran status, and log wages as the outcome.

\end{frame}


\begin{frame}{A harmless coding choice}
Suppose income $Y_i$ is measured in dollars. If the analysis is re-run in cents: $Y_i' = 100\,Y_i$.

\vspace{0.2cm}
If you log-transform first (standard practice in applied micro) the two are related by a constant shift:
\[
\log(100\,Y_i) \;=\; \log(Y_i) + \log(100).
\]

\vspace{0.2cm}
So switching currency units is, after logging, exactly an additive shift of the outcome by a constant $c = \log(100)$.

\vspace{0.3cm}
\centering \alert{Should this change your estimated causal effect?}
\end{frame}


\begin{frame}{A coding choice that should not matter}

\vspace{0.2cm}
\citet{sloczynski2025abadie} report exactly this comparison, holding the data and specification fixed:

\vspace{0.2cm}
\begin{itemize}
\item Their \textbf{normalized} weighting estimators ($\hat\tau_u$, $\hat\tau_{a,10}$) and 2SLS give (essentially) the \textbf{same} estimate whether wages are coded in cents or dollars.
\item Several \textbf{unnormalized} weighting estimators ($\hat\tau_a$, $\hat\tau_t$, $\hat\tau_{a,0}$) give \textbf{different} estimates, in one specification, even flipping from negative to positive.
\end{itemize}

\vspace{0.3cm}
\centering \alert{Same data, same question but different conclusion, depending only on currency units.}
\end{frame}

\begin{frame}{What explains the coding sensitivity?}
The cents-versus-dollars example suggests the problem is not the data or the IV design, but how the estimator combines outcomes.

\vspace{0.25cm}
Write any such estimator as a weighted sum, $\hat\tau = \sum_{i=1}^N \omega_i Y_i$. Shifting every outcome by a constant $c$ gives
\[
\hat\tau(Y+c) = \hat\tau(Y) + c\sum_{i=1}^N \omega_i
\quad\Longrightarrow\quad
\hat\tau(Y+c)=\hat\tau(Y) \iff \sum_{i=1}^N \omega_i = 0.
\]

\vspace{0.3cm}
\centering
\alert{Translation invariance is a finite-sample property of the outcome weights $\omega_i$ not of the data or the IV design.}
\end{frame}

\begin{frame}{A broader diagnostic lens}
This weighted-sum representation is the entry point to a broader framework developed by \citet{knaus2024treatment}: once $\omega_i$ is known, it can audit an estimator well beyond invariance:

\begin{itemize}
\item \textbf{Effective sample size}: how much of the data is really used
\item \textbf{Negative \& extreme weights} : whether a few observations drive the estimate
\item \textbf{Covariate balance}: whether the comparison resembles a credible experiment
\end{itemize}

\vspace{0.25cm}
Knaus derives these weights explicitly for DML and generalized random forest estimators.

\vspace{0.2cm}
\centering
\alert{My thesis derives the same closed-form weights for kappa-based LATE estimators, extending this diagnostic lens beyond translation invariance.}
\end{frame}


\begin{frame}{Research question}
The previous example motivates a broader question:

\vspace{0.2cm}
\begin{block}{Thesis research question}
Can outcome weights make classical kappa-weighting and modern causal machine learning estimators of the LATE more comparable, interpretable, and diagnostically transparent?
\end{block}

\vspace{0.25cm}
This thesis connects two ideas:

\begin{itemize}
\item \textbf{Translation invariance:} the total weight sum $\sum_i \omega_i$ determines whether an estimator is stable to harmless outcome shifts.
\item \textbf{Estimator interpretability:} the distribution of $\omega_i$ reveals which observations drive the estimate, with which sign, and how strongly.
\end{itemize}



\vspace{0.2cm}
\centering
\alert{The goal is to move beyond the point estimate and audit how the estimator uses the sample.}
\end{frame}






% ============================================================
\section{Empirical Approach}
% ============================================================

\begin{frame}{Empirical approach in my thesis}
\textbf{Three famous IV applications from the literature}:
\begin{enumerate}
\item \citet{angrist1990lifetime}: Vietnam draft lottery as an instrument for veteran status
\item \citet{card1993using}: college proximity as an instrument for education
\item \citet{angrist1996children}: sibling sex composition as an instrument for fertility & labor supply
\end{enumerate}

\vspace{0.3cm}
For each: compute outcome weights for all five kappa estimators (plus DML/GRF benchmarks using alternative nuisance learners), then assess ESS, covariate balance, and translation invariance.
\vspace{0.2cm}

\alert{This makes classical weighting estimators and causal ML estimators comparable through the same finite-sample objects.}
\end{frame}



\begin{frame}{Emerging patterns}

The outcome-weight lens reveals more than translation invariance:

\vspace{0.2cm}

\begin{itemize}
\item \textbf{Design dominates learner:} in clean IV designs, kappa and DML estimators can produce very similar weights and balance diagnostics.

\vspace{0.1cm}

\item \textbf{Weak support becomes visible:} limited overlap or weak first stages show up as low ESS, extreme weights, and unstable estimates.

\vspace{0.1cm}

\item \textbf{Causal ML needs auditing:} different nuisance learners can imply different outcome weights, even when point estimates look similar.
\end{itemize}

\vspace{0.3cm}
\centering
\alert{This suggests a broader PhD direction: making causal ML estimators flexible, but also interpretable and auditable.}

\end{frame}





% ============================================================
\section{Outlook}
% ============================================================
\begin{frame}{From thesis to PhD agenda}
My thesis is a proof of concept: outcome weights can make LATE estimators more transparent in finite samples.

\vspace{0.25cm}

\begin{block}{Possible PhD direction}
When do flexible causal ML estimators improve, replicate, or destabilize classical causal estimates --- and can diagnostics reveal this before we trust the point estimate?
\end{block}

\vspace{0.25cm}

Three directions I would be excited to explore:

\begin{itemize}
\item \textbf{Heterogeneous effects:} audit CATE / CLATE estimates locally, not only one average LATE.

\item \textbf{Rich nuisance learners:} study how XGBoost, deep nets, or LLM/text embeddings inside DML affect the final causal estimate.

\item \textbf{Estimator explainability:} explain causal estimates through outcome weights, influence, balance, and effective sample size.
\end{itemize}

\vspace{0.25cm}
\centering
\alert{The long-run goal is estimator-level explainability for causal AI systems.}

\end{frame}


\begin{frame}{Two concrete paths}

\textbf{1. Local support for heterogeneous causal effects}

\vspace{0.1cm}
CATE / CLATE estimators can produce effect estimates for specific covariate regions.  
But a local estimate is only useful if it is supported by enough comparable observations.

\vspace{0.15cm}


\vspace{0.3cm}

\textbf{2. Text and LLMs as nuisance information in DML}

\vspace{0.1cm}
In many applications, confounders are not only tabular variables but text:
clinical notes, reviews, news articles, job descriptions, policy documents.

\vspace{0.15cm}
LLM embeddings could enter DML as high-dimensional nuisance information, but then the key question becomes:

\vspace{0.1cm}


\vspace{0.25cm}
\centering
\alert{Outcome-weight diagnostics could help audit this distinction.}

\end{frame}


\bibliographystyle{apalike}
\bibliography{references}

\end{document}
